%%% This file contains definitions of various useful macros and environments %%%
%%% Please add more macros here instead of cluttering other files with them. %%%

%%% Switches based on thesis type

\def\TypeBc{bc}
\def\TypeMgr{mgr}
\def\TypePhD{phd}
\def\TypeRig{rig}

\def\ThesisTypeName{doctoral}
\def\ThesisTypeTitle{DOCTORAL THESIS}
\def\ThesisTypeCSGenitive{disertační}

%%% Switches based on study program language

\def\LangCS{cs}
\def\LangEN{en}

\ifx\StudyLanguage\LangCS
\else\ifx\StudyLanguage\LangEN
\else\PackageError{thesis}{Unknown study language.}{Please check the definition of StudyLanguage in metadata.tex.}
\fi\fi

%%% Minor tweaks of style

% These macros employ a little dirty trick to convince LaTeX to typeset
% chapter headings sanely, without lots of empty space above them.
% Feel free to ignore.
\makeatletter
\def\@makechapterhead#1{
  {\parindent \z@ \raggedright \normalfont
   \Huge\bfseries \thechapter\quad #1
   \par\nobreak
   \vskip 20\p@
}}
\def\@makeschapterhead#1{
  {\parindent \z@ \raggedright \normalfont
   \Huge\bfseries #1
   \par\nobreak
   \vskip 20\p@
}}
\makeatother

% This macro defines a chapter, which is not numbered, but is included
% in the table of contents.
\def\chapwithtoc#1{
\chapter*{#1}
\addcontentsline{toc}{chapter}{#1}
}

\def\secwithtoc#1{
\section*{#1}
\addcontentsline{toc}{section}{#1}
}

% Draw black "slugs" whenever a line overflows, so that we can spot it easily.
\overfullrule=5mm
\emergencystretch=1em

\AtBeginDocument{
\hyphenation{SOM-Hun-ter}
\hyphenation{Wi-de-spread}
}

%%% Macros for definitions, theorems, claims, examples, ... (requires amsthm package)

\theoremstyle{plain}
\newtheorem{thm}{Theorem}
\newtheorem{lemma}[thm]{Lemma}
\newtheorem{claim}[thm]{Claim}

\theoremstyle{plain}
\newtheorem{defn}{Definition}
\newtheorem{problem}{Problem}

\theoremstyle{remark}
\newtheorem*{cor}{Corollary}
\newtheorem*{rem}{Remark}
\newtheorem*{example}{Example}

%%% Style of captions of floating objects (figures etc.)

\ifcsname DeclareCaptionStyle\endcsname
\DeclareCaptionStyle{thesis}{style=base,font=small,labelfont=bf,labelsep=quad}
\captionsetup{style=thesis}
\captionsetup[algorithm]{style=thesis,singlelinecheck=off}
\captionsetup[listing]{style=thesis,singlelinecheck=off}
\fi

%%% An environment for typesetting of program code and input/output
%%% of programs.
\DefineVerbatimEnvironment{code}{Verbatim}{fontsize=\small, frame=single}

% Settings for lstlisting -- program listing with syntax highlighting
\ifcsname lstset\endcsname
\lstdefinestyle{C++}{
  language=C++,
  tabsize=2,
  upquote=true,
  % captions at the bottom
  captionpos=b,
  % number captions like "Listing 1., ..."
  numberbychapter=false,
  numbers=left,
  numberstyle=\tiny\color{black!50},
  backgroundcolor=\color{black!3},
  showstringspaces=false,
  basicstyle=\footnotesize\ttfamily\color{black!75},
  identifierstyle=\bfseries\color{black},
  commentstyle=\color{green!50!black},
  stringstyle=\color{red!75!black},
  keywordstyle=\color{blue!75!black},
  morekeywords={concept,requires,co_await,co_yield,co_return,consteval,constexpr,decltype,final,noexcept,override,static_assert,thread_local},
  classoffset=1,
  morekeywords={noarr,cellato,std,Kokkos,RAJA},
  keywordstyle=\color{teal!75!black},
  classoffset=0,
  literate=
  {__0}{{\textcolor{teal!25!purple}{0}}}1
  {__1}{{\textcolor{teal!25!purple}{1}}}1
  {__2}{{\textcolor{teal!25!purple}{2}}}1
  {__3}{{\textcolor{teal!25!purple}{3}}}1
  {__4}{{\textcolor{teal!25!purple}{4}}}1
  {__5}{{\textcolor{teal!25!purple}{5}}}1
  {__6}{{\textcolor{teal!25!purple}{6}}}1
  {__7}{{\textcolor{teal!25!purple}{7}}}1
  {__8}{{\textcolor{teal!25!purple}{8}}}1
  {__9}{{\textcolor{teal!25!purple}{9}}}1
  {__.}{{\textcolor{teal!25!purple}{.}}}1
  {__f}{{\textcolor{teal!25!purple}{f}}}1
}
\lstset{style=C++}

\fi

%%% The field of all real and natural numbers
\newcommand{\R}{\mathbb{R}}
\newcommand{\N}{\mathbb{N}}
\newcommand{\Asymp}[1]{\ensuremath{\mathcal{O}\left(#1\right)}}

%%% Nicer "C++" in text
\newcommand{\cpp}{C\texttt{++}\xspace}
\newcommand{\cppeleven}{C\texttt{++}11\xspace}
\newcommand{\cppseventeen}{C\texttt{++}17\xspace}
\newcommand{\cpptwenty}{C\texttt{++}20\xspace}
\newcommand{\cpptwentythree}{C\texttt{++}23\xspace}
\newcommand{\cpptwentysix}{C\texttt{++}26\xspace}

%%% Useful operators for statistics and probability
\DeclareMathOperator{\pr}{\textsf{P}}
\DeclareMathOperator{\E}{\textsf{E}\,}
\DeclareMathOperator{\var}{\textrm{var}}
\DeclareMathOperator{\sd}{\textrm{sd}}
\DeclareMathOperator{\subs}{\textsc{Sub}}

% contributions
\newcounter{contribution}
\def\contribution#1{
\refstepcounter{contribution}
\chapter*{Contribution \arabic{contribution}\\#1}
\addcontentsline{toc}{chapter}{Contribution \arabic{contribution}\\ #1}
}

\crefname{contribution}{contribution}{contributions}
\Crefname{contribution}{Contribution}{Contributions}

%%% TODO items: remove before submitting :)
\newcommand{\xxx}[1]{\textcolor{red!}{#1}}

\newenvironment{XXX}
  {\par\color{red}}
  {\par}


\setlength{\marginparpush}{0pt}
\setlength{\marginparsep}{2em}

\newcommand{\contribmark}[1]{%
\hyperref[#1]{\begin{tikzpicture}[ultra thick, font=\sf\bfseries]
\node[regular polygon, regular polygon sides=3, rounded corners=0.3ex, inner sep=.3ex,draw] (ll) {\ref*{#1}};
\end{tikzpicture}}}
\newcommand{\margincontrib}[1]{\marginnote{\contribmark{#1}}}

\newcommand{\publishedas}[1]{\begin{description}\item[Published as]\begin{refsection}\fullcite{#1}\end{refsection}\end{description}}

% ONLINE VERSION without the manuscripts
\newif\ifonlineversion
\onlineversiontrue
%\onlineversionfalse
\newcommand\onlineomit[1]{\ifonlineversion\vfill\begin{center}\fbox{\parbox{16.5em}{\footnotesize A re-printed manuscript that would appear here is omitted in the electronic version of the thesis.}}\end{center}\vfill\else#1\fi}

% publication contributions
\def\contF{\begin{tikzpicture}\fill circle(0.666ex);\end{tikzpicture}\xspace}
\def\contP{\begin{tikzpicture}\begin{scope}[rotate=90]\clip circle(0.666ex);\fill (-0.666ex,-0.666ex) rectangle(-0.1ex,0.666ex);\end{scope}\draw circle(0.666ex);\end{tikzpicture}\xspace}
\def\contN{\begin{tikzpicture}\draw[densely dotted] circle(0.666ex);\end{tikzpicture}\xspace}
%%% Detailed settings of bibliography

\ifx\citet\undefined\else

% Maximum number of authors of a single work. If exceeded, "et al." is used.
%\ExecuteBibliographyOptions{maxnames=2}
% The same setting specific to citations using \citet{...}
\ExecuteBibliographyOptions{maxcitenames=2}
% The same settings specific to the list of literature
%\ExecuteBibliographyOptions{maxbibnames=2}

% Shortening first names of authors: "E. A. Poe" instead of "Edgar Allan Poe"
%\ExecuteBibliographyOptions{giveninits}
% The same without dots ("EA Poe")
%\ExecuteBibliographyOptions{terseinits}

% If your bibliography entries are hard to break into lines, try this mode:
%\ExecuteBibliographyOptions{block=ragged}

% Possibly reverse the names of the authors with the non-ISO styles:
%\DeclareNameAlias{default}{family-given}

% Use caps-and-small-caps for family names in ISO 690 style.
\let\familynameformat=\textsc

% We want to separate multiple authors in citations by commas
% (while we use semicolons in the bibliography as per the ISO standard)
\DeclareDelimFormat[textcite]{multinamedelim}{\addcomma\space}
\DeclareDelimFormat[textcite]{finalnamedelim}{\space and~}

\fi

% latexdiff preamble commands:
%DIF UNDERLINE PREAMBLE %DIF PREAMBLE
\RequirePackage[normalem]{ulem} %DIF PREAMBLE
\RequirePackage{color}\definecolor{RED}{rgb}{1,0,0}\definecolor{BLUE}{rgb}{0,0,1} %DIF PREAMBLE
\providecommand{\DIFadd}[1]{{\protect\color{green!50!black}\uwave{#1}}} %DIF PREAMBLE
\providecommand{\DIFdel}[1]{{\protect\color{red}\sout{#1}}}                      %DIF PREAMBLE
%DIF SAFE PREAMBLE %DIF PREAMBLE
\providecommand{\DIFaddbegin}{} %DIF PREAMBLE
\providecommand{\DIFaddend}{} %DIF PREAMBLE
\providecommand{\DIFdelbegin}{} %DIF PREAMBLE
\providecommand{\DIFdelend}{} %DIF PREAMBLE
\providecommand{\DIFmodbegin}{} %DIF PREAMBLE
\providecommand{\DIFmodend}{} %DIF PREAMBLE
%DIF FLOATSAFE PREAMBLE %DIF PREAMBLE
\providecommand{\DIFaddFL}[1]{\DIFadd{#1}} %DIF PREAMBLE
\providecommand{\DIFdelFL}[1]{\DIFdel{#1}} %DIF PREAMBLE
\providecommand{\DIFaddbeginFL}{} %DIF PREAMBLE
\providecommand{\DIFaddendFL}{} %DIF PREAMBLE
\providecommand{\DIFdelbeginFL}{} %DIF PREAMBLE
\providecommand{\DIFdelendFL}{} %DIF PREAMBLE
